\chapter{2011年天津卷（理）}

\section{单选题}

\begin{problem}
i 是虚数单位，复数 \( \frac{1-3\i}{1-\i}= \)（\quad）

\choices
    {\(2 + \i\)}
    {\(2 - \i\)}
    {\(-1 + 2\i\)}
    {\(-1 - 2\i\)}
\end{problem}

\begin{problem}
设 \(x\)，\(y \in \R\)，则“\(x \geqslant 2\) 且 \(y \geqslant 2\)”是“\(x^2 + y^2 \geqslant 4\)”的（\quad）

\choices
    {充分而不必要条件}
    {必要而不充分条件}
    {充分必要条件}
    {既不充分也不必要条件}
\end{problem}

\begin{problem}
阅读下面的程序框图，运行相应的程序，则输出 \(i\) 的值为（\quad）

\texfigure[max width=0.70\linewidth]{img_repaint/2011/tianjin_science/q3_fig1.tex}

\choices
    {3}
    {4}
    {5}
    {6}
\end{problem}

\begin{problem}
已知 \( \{a_{n}\} \) 为等差数列，其公差为 \(-2\)，且 \( a_{7} \) 是 \( a_{3} \) 与 \( a_{9} \) 的等比中项，\( S_{n} \) 为 \( \{a_{n}\} \) 的前 n 项和，\( n \in \N^{*} \)，则 \(S_{10}\) 的值为（\quad）

\choices
    {\(-110\)}
    {\(-90\)}
    {90}
    {110}
\end{problem}

\begin{problem}
在 \(\left(\frac{\sqrt{x}}{2} - \frac{2}{\sqrt{x}}\right)^{6}\) 的二项展开式中，\(x^{2}\) 的系数为（\quad）

\choices
    {\(-\frac{15}{4}\)}
    {\( \frac{15}{4} \)}
    {\(-\frac{3}{8}\)}
    {\(\frac{3}{8}\)}
\end{problem}

\begin{problem}
如图所示，在 \(\triangle ABC\) 中，\(D\) 是边 \(AC\) 上的点，且 \(AB = AD\)，\(2AB = \sqrt{3}BD\)，\(BC = 2BD\)，则 \(\sin C\) 的值为（\quad）

\bitmapfigure[max width=0.34\linewidth]{img/2011/tianjin_science/q6_fig1.png}

\choices
    {\(\frac{\sqrt{3}}{3}\)}
    {\(\frac{\sqrt{3}}{6}\)}
    {\(\frac{\sqrt{6}}{3}\)}
    {\(\frac{\sqrt{6}}{6}\)}
\end{problem}

\begin{problem}
已知 \(a = 5^{\log_{2} 3^{4}}\)，\(b = 5^{\log_{4} 3^{6}}\)，\(c = \left(\frac{1}{5}\right)^{\log_{3} 0.3}\)，则（\quad）

\choices
    {\(a > b > c\)}
    {\(b > a > c\)}
    {\(a > c > b\)}
    {\(c > a > b\)}
\end{problem}

\begin{problem}
对实数 \(a\) 与 \(b\)，定义运算“\(\otimes\)”： \(a \otimes b = \begin{cases} a, & a - b \leqslant 1, \\ b, & a - b > 1. \end{cases}\) 设函数 \(f(x) = (x^2 - 2) \otimes (x - x^2)\)，\(x \in \R\). 若函数 \(y = f(x) - c\) 的图象与 \(x\) 轴恰有两个公共点，则实数 \(c\) 的取值范围是（\quad）

\choices
    {\((- \infty, -2] \cup \left(-1, \frac{3}{2}\right)\)}
    {\((-\infty, -2] \cup \left(-1, -\frac{3}{4}\right)\)}
    {\(\left(-1, \frac{1}{4}\right) \cup \left(\frac{1}{4}, +\infty\right)\)}
    {\(\left(-1, -\frac{3}{4}\right) \cup \left[\frac{1}{4}, +\infty\right)\)}
\end{problem}

\section{填空题}

\begin{problem}
一支田径队有男运动员 48 人，女运动员 36 人，若用分层抽样的方法从该队的全体运动员中抽取一个容量为 21 的样本，则抽取男运动员的人数为\fillinblank{}.
\end{problem}

\begin{problem}
一个几何体的三视图如图所示（单位：m），则这个几何体的体积为\fillinblank{} \(\text{m}^{3}\).

\bitmapfigure[max width=0.56\linewidth]{img/2011/tianjin_science/q10_fig1.png}

\end{problem}

\begin{problem}
已知抛物线 C 的参数方程为 \( \left\{\begin{aligned} x &= 8t^{2}, \\ y &= 8t, \end{aligned}\right. \) （\(t\) 为参数），若斜率为 1 的直线经过抛物线 \(C\) 的焦点，且与圆 \( (x - 4)^{2} + y^{2} = r^{2} \) （\(r > 0\)）相切，则 \(r =\)\fillinblank{}.
\end{problem}

\begin{problem}
如图，已知圆中两条弦 \(AB\) 与 \(CD\) 相交于点 \(F\)，\(E\) 是 \(AB\) 延长线上一点，且 \(DF = CF = \sqrt{2}\)，\(AF:FB:BE = 4:2:1\).
若 \(CE\) 与圆相切，则 \(CE\) 的长为\fillinblank{}.

\bitmapfigure[max width=0.42\linewidth]{img/2011/tianjin_science/q12_fig1.png}

\end{problem}

\begin{problem}
已知集合 \(A = \{x \in \R \mid |x + 3| + |x - 4| \leqslant 9\}\)，\(B = \left\{x \in \R \mid x = 4t + \frac{1}{t} - 6, t \in (0, +\infty)\right\}\)，则集合 \(A \cap B = \)\fillinblank{}.
\end{problem}

\begin{problem}
已知直角梯形 \(ABCD\) 中，\(AD \parallel BC\)，\(\angle ADC = 90^{\circ}\)，\(AD = 2\)，\(BC = 1\)，\(P\) 是腰 \(DC\) 上的动点，则 \(\left|\overrightarrow{PA} + 3\overrightarrow{PB}\right|\) 的最小值为\fillinblank{}.
\end{problem}

\section{解答题}

\begin{problem}
已知函数 \( f(x) = \tan \left(2x + \frac{\pi}{4}\right) \).
\begin{enumerate}
\item 求 \( f(x) \) 的定义域与最小正周期；
\item 设 \(\alpha \in (0, \frac{\pi}{4})\)，若 \(f\left(\frac{\alpha}{2}\right) = 2\cos 2\alpha\)，求 \(\alpha\) 的大小.
\end{enumerate}
\end{problem}

\begin{problem}
学校游园活动有这样一个游戏项目：甲箱子里装有3个白球，2个黑球，乙箱子里装有1个白球，2个黑球，这些球除颜色外完全相同. 每次游戏从这两个箱子里各随机摸出2个球，若摸出的白球不少于2个，则获奖. （每次游戏结束后将球放回原箱）
\begin{enumerate}
\item 求在一次游戏中
\begin{enumerate}
    \item 摸出 3 个白球的概率；
    \item 获奖的概率.
\end{enumerate}
\item 求在两次中获奖次数 X 的分布列及数学期望 \( E(X) \).
\end{enumerate}
\end{problem}

\begin{problem}
如图，在三棱柱 \(ABC - A_{1}B_{1}C_{1}\) 中，\(H\) 是正方形 \(AA_{1}B_{1}B\) 的中心，\(AA_{1} = 2\sqrt{2}\)，\(C_{1}H \perp\) 平面 \(AA_{1}B_{1}B\)，且 \(C_{1}H = \sqrt{5}\).
\begin{enumerate}
\item 求异面直线 \(AC\) 与 \(A_{1}B_{1}\) 所成角的余弦值；
\item 求二面角 \(A - A_{1}C_{1} - B_{1}\) 的正弦值；
\item 设 \(N\) 为棱 \(B_{1}C_{1}\) 的中点，点 \(M\) 在平面 \(AA_{1}B_{1}B\) 内，且 \(MN \perp\) 平面 \(A_{1}B_{1}C_{1}\)，求线段 \(BM\) 的长.
\end{enumerate}
\bitmapfigure[max width=0.42\linewidth]{img/2011/tianjin_science/q17_fig1.png}
\end{problem}

\begin{problem}
已知 \(a > 0\)，函数 \(f(x) = \ln x - ax^2\)，\(x > 0\).（\(f(x)\) 的图象连续不断）
\begin{enumerate}
\item 求 \( f(x) \) 的单调区间；
\item 当 \(a = \frac{1}{8}\) 时，证明：存在 \(x_0 \in (2, +\infty)\)，使 \(f(x_0) = f\left(\frac{3}{2}\right)\).
\item 若存在均属于区间 \([1,3]\) 的 \(\alpha\)，\(\beta\)，且 \(\beta - \alpha \geqslant 1\)，使 \(f(\alpha) = f(\beta)\). 证明： \(\frac{\ln 3 - \ln 2}{5} \leqslant a \leqslant \frac{\ln 2}{3}\).
\end{enumerate}
\end{problem}

\begin{problem}
已知数列 \(\{a_{n}\}\) 与 \(\{b_{n}\}\) 满足： \(b_{n}a_{n} + a_{n + 1} + b_{n + 1}a_{n + 2} = 0\)，\(b_{n} = \frac{3 + (-1)^{n}}{2}\)，\(n \in \N^{*}\)，且 \(a_{1} = 2\)，\(a_{2} = 4\)，
\begin{enumerate}
\item 求 \(a_{3}\)，\(a_{4}\)，\(a_{5}\) 的值；
\item 设 \(c_{n} = a_{2n - 1} + a_{2n + 1}\)，\(n \in \N^{*}\)，证明： \(\{c_{n}\}\) 是等比数列；
\item 设 \(S_{k} = a_{2} + a_{4} + \dots + a_{2k}\)，\(k \in \N^{*}\)，证明： \(\sum_{k=1}^{n} \frac{S_{k}}{a_{k}} < \frac{7}{6} (n \in \N^{*})\).
\end{enumerate}
\end{problem}

\begin{problem}
在平面直角坐标系 \(xOy\) 中，点 \(P(a,b)\) （\(a > b > 0\)） 为动点，\(F_{1}\)，\(F_{2}\) 分别为椭圆 \(\frac{x^2}{a^2} + \frac{y^2}{b^2} = 1\) 的左，右焦点. 已知 \(\triangle F_1PF_2\) 为等腰三角形.
\begin{enumerate}
\item 求椭圆的离心率 \(e\)；
\item 设直线 \(PF_{2}\) 与椭圆相交于 \(A\)，\(B\) 两点，\(M\) 是直线 \(PF_{2}\) 上的点，满足 \(\overrightarrow{AM} \cdot \overrightarrow{BM} = -2\)，求点 \(M\) 的轨迹方程.
\end{enumerate}
\end{problem}
